
  

\fchapter{Ej 11 Pág 50}
\fsection[\textcolor{waveBlue2}{\boldit{Pensamiento crítico}}]{\boldit{
		\textcolor{waveBlue2}{UTILIZA LAS TIC.} Busca información en internet sobre las zonas de bajas
		emisiones y responde a las siguientes preguntas:
}} 

\begin{solution}
	
\begin{enumerate}[label=\alph*)]
	\item{\boldit{¿Qué se pretende con la creación de las zonas de bajas emisiones?}}
		\vspace{1em}
		La creación de las zonas de bajas emisiones tiene como objetivo principal reducir la contaminación del aire y mejorar la calidad de vida de la población. Con estas zonas se busca disminuir las emisiones de gases contaminantes producidas por el tráfico, fomentar el uso de medios de transporte más sostenibles y contribuir a la lucha contra el cambio climático, además de proteger la salud de las personas.
		\vspace{1em}

	\item{\boldit{¿Qué características debe tener una población para ser considerada zona de bajas emisiones?}}

		\vspace{1em}
		Una población debe contar con un elevado nivel de tráfico y problemas de contaminación atmosférica para ser considerada zona de bajas emisiones. Normalmente se trata de municipios con una alta densidad de población, donde se superan o se corre el riesgo de superar los límites de contaminación establecidos por la normativa. Además, deben disponer de infraestructuras y medidas que permitan controlar el acceso de los vehículos y promover alternativas de movilidad sostenible.
		\vspace{1em}

	\item{\boldit{¿Qué vehiculos pueden circular por las zonas de bajas emisiones?}}

		\vspace{1em}
		En las zonas de bajas emisiones pueden circular los vehículos menos contaminantes, como los eléctricos, híbridos enchufables y aquellos que cumplen las normativas de emisiones más recientes. Generalmente se permite el acceso a los vehículos con distintivos ambientales favorables, mientras que los más antiguos y contaminantes tienen restricciones o directamente no pueden acceder, salvo algunas excepciones como servicios de emergencia o transporte público.
		\vspace{1em}

\end{enumerate}

\end{solution}
